% État de l'art, limité à ce qui éclaire les choix du projet.

\chapter{État de l'art}
\label{chap:etat-art}

\section{Modèles de buts}

\paragraph{Poisson par équipe.} Depuis Maher~\cite{maher1982}, on modélise les buts de chaque équipe
par une loi de Poisson dont le logarithme de la moyenne combine l'attaque de l'équipe, la défense de
l'adversaire et l'avantage du terrain. Les deux comptages sont supposés indépendants sachant ces
forces ; le total suit alors une loi de Poisson de paramètre \(\lambda_{\mathrm{dom}} +
\lambda_{\mathrm{ext}}\).

\paragraph{Dépendance entre équipes.} Dixon et Coles~\cite{dixon1997} corrigent les quatre scores les
plus faibles (0-0, 1-0, 0-1, 1-1) par un paramètre \(\rho\), et pondèrent les matchs anciens par une
décroissance exponentielle. Karlis et Ntzoufras~\cite{karlis2003} proposent un Poisson bivarié, avec
une composante commune aux deux équipes. Ces raffinements visent surtout le score exact ; le projet
mesure s'ils changent la loi du \emph{total} (section~\ref{sec:resultats-modeles}).

\paragraph{Force des équipes.} Le classement Elo~\cite{elo1978}, conçu pour les échecs, met à jour une
note après chaque match selon l'écart entre résultat et attente. Hvattum et Arntzen~\cite{hvattum2010}
montrent qu'une note Elo est une variable efficace pour prédire les résultats en football. Faute de
source gratuite (le service ClubElo est passé derrière une authentification), le projet calcule son
propre Elo.

\paragraph{Dispersion.} La loi de Poisson impose une variance égale à la moyenne. Le test de Cameron
et Trivedi~\cite{cameron1990} détecte une surdispersion, qui justifierait une binomiale négative.

\section{Évaluer une prédiction probabiliste}

Une règle de score \emph{strictement propre} n'est minimisée en espérance que par la vraie loi
\cite{gneiting2007} : elle récompense l'honnêteté de la probabilité annoncée. Le log-loss en est une,
sensible à toute la loi. Le RPS (\emph{ranked probability score}) tient compte de l'ordre des
catégories, ce qui convient aux résultats sportifs~\cite{constantinou2012}. Pour une loi discrète, le
PIT randomisé~\cite{czado2009} vérifie la calibration d'ensemble. Comparer deux modèles sur les mêmes
matchs se fait par la différence de leurs pertes ; le test de Diebold et Mariano~\cite{diebold1995}
en donne la significativité, et le bootstrap par blocs~\cite{kunsch1989} un intervalle de confiance
qui respecte la dépendance entre matchs d'une même journée.

\section{Le marché des paris}

Les cotes agrègent l'information des bookmakers et des parieurs ; Dixon et Coles s'en servaient déjà
pour tester leur modèle~\cite{dixon1997}. Le projet les utilise comme \emph{référence}, jamais comme
variable : un modèle qui lirait les cotes ne prédirait plus le football, mais le marché (ADR-0036).
